\section{实验四：相机标定}

\subsection{实验目的}

本实验目标是整理相机内参、畸变参数以及相机与工作平面之间位姿关系的计算流程，为后续从图像像素坐标推算机械臂抓取坐标提供程序基础。本组未通过真实机械臂运动检验转换后的抓取坐标。

\subsection{相机成像模型}

实验采用针孔相机模型。空间点 $P_C=[X_C,Y_C,Z_C]^T$ 在相机坐标系下投影到像素平面：
\[
s
\begin{bmatrix}
u\\v\\1
\end{bmatrix}
=
K
\begin{bmatrix}
X_C\\Y_C\\Z_C
\end{bmatrix},
\quad
K=
\begin{bmatrix}
f_x&0&c_x\\
0&f_y&c_y\\
0&0&1
\end{bmatrix}.
\]
其中 $K$ 为相机内参矩阵，$f_x,f_y$ 为像素焦距，$c_x,c_y$ 为主点坐标。实际镜头还存在径向畸变和切向畸变，需要通过标定估计畸变系数。

\subsection{内参与畸变标定}

内参标定程序为 \path{daran/task3/camera_calibration.py}。程序支持实时采集和从已有图片标定两种方式。实验使用 $7\times7$ 内角点棋盘格，方格边长为 \SI{0.008}{m}，有效标定图像共 30 张。程序流程如下：

\begin{enumerate}
  \item 使用 \texttt{cv2.findChessboardCorners} 检测棋盘格角点。
  \item 使用 \texttt{cv2.cornerSubPix} 对角点进行亚像素优化。
  \item 构造棋盘格世界坐标点。
  \item 调用 \texttt{cv2.calibrateCamera} 求解内参和畸变。
  \item 保存 \path{camera_calibration.json}、\path{camera_calibration.npz} 和去畸变示例图。
\end{enumerate}

标定结果如表 \ref{tab:camera-calib} 所示。

\begin{table}[H]
  \centering
  \caption{相机内参标定结果}
  \label{tab:camera-calib}
  \small
  \begin{tabularx}{0.72\textwidth}{L{3.5cm}Y}
    \toprule
    项目 & 数值 \\
    \midrule
    图像尺寸 & $640\times480$ \\
    有效图像数 & 30 \\
    RMS 重投影误差 & 0.2025 px \\
    平均重投影误差 & 0.2025 px \\
    $f_x$ & 381.2781 \\
    $f_y$ & 508.3521 \\
    $c_x$ & 319.2541 \\
    $c_y$ & 239.6599 \\
    \bottomrule
  \end{tabularx}
\end{table}

标定得到的相机矩阵为：
\[
K=
\begin{bmatrix}
381.2781 & 0 & 319.2541\\
0 & 508.3521 & 239.6599\\
0 & 0 & 1
\end{bmatrix}.
\]
畸变系数采用 OpenCV 默认顺序，可写为：
\[
(k_1,k_2,p_1,p_2,k_3)
=(0.1455,-0.2150,-0.0025,0.0029,0.0715).
\]
其中 $k_1,k_2,k_3$ 描述径向畸变，$p_1,p_2$ 描述切向畸变。径向畸变主要与镜头边缘区域的弯曲有关，切向畸变则与镜头和成像平面不完全平行有关。去畸变后程序还计算了新的相机矩阵
\[
K_{\text{new}}=
\begin{bmatrix}
383.3236 & 0 & 322.0648\\
0 & 509.6299 & 237.8273\\
0 & 0 & 1
\end{bmatrix},
\]
对应有效 ROI 为 $(6,6,628,466)$。这意味着去畸变图像会裁掉少量边缘区域，以减少无效黑边对后续识别的影响。

重投影误差是评价标定质量的重要指标。标定完成后，程序会将棋盘格世界坐标点通过估计得到的相机参数重新投影到图像平面，再与实际检测到的角点位置比较。若误差较小，说明内参、畸变参数和各标定图像的位姿能够较好解释观测数据。本实验平均重投影误差约为 0.2025 px，表明角点检测和内参求解具有较好一致性。

标定图像的采集质量会直接影响结果。理想情况下，棋盘格应覆盖图像中心和边缘区域，并包含不同倾角和不同距离的姿态；若图像都集中在画面中央，畸变参数约束不足；若棋盘格角点模糊、反光或遮挡，亚像素角点会产生偏差。因此相机标定不仅是一次函数调用，更依赖采集过程中的姿态多样性和图像清晰度。

标定程序中棋盘格三维点生成、角点检测和相机标定调用如代码 \ref{lst:exp4-code} 所示。该部分是后续所有像素坐标转换的基础。

\lstinputlisting[
  language=Python,
  caption={棋盘格角点检测与相机内参标定核心代码},
  label={lst:exp4-code}
]{code/exp4_camera_calibration.py}

\begin{figure}[H]
  \centering
  \includegraphics[width=0.72\textwidth]{task4_undistorted_example.jpg}
  \caption{相机去畸变示例图}
  \label{fig:task4-undistorted}
\end{figure}

\subsection{工作空间坐标转换}

外参估计程序为 \path{daran/task3/pose_projection.py}。在工作平面上放置棋盘格后，程序通过 \texttt{cv2.solvePnP} 求解世界坐标系到相机坐标系的变换矩阵，并保存 \path{camera_extrinsic.json}。实验得到的外参重投影误差约为 0.0215 px。

本实验中棋盘格坐标系到相机坐标系的平移向量为：
\[
t=
\begin{bmatrix}
-0.0158 & -0.0316 & 0.1572
\end{bmatrix}^T\text{m},
\]
说明棋盘格原点相对于相机光心位于约 \SI{15.7}{cm} 的深度处。旋转向量为：
\[
r=
\begin{bmatrix}
-0.0234 & -0.0160 & 0.0074
\end{bmatrix}^T,
\]
其数值较小，说明棋盘格平面相对于相机坐标系的倾斜角不大。对应的世界到相机齐次变换矩阵为：
\[
T_{C}^{W}=
\begin{bmatrix}
0.9998 & -0.0072 & -0.0161 & -0.0158\\
0.0076 & 0.9997 & 0.0233 & -0.0316\\
0.0159 & -0.0234 & 0.9996 & 0.1572\\
0 & 0 & 0 & 1
\end{bmatrix}.
\]
该矩阵的前三列描述棋盘格坐标轴在相机坐标系中的方向，最后一列描述棋盘格原点在相机坐标系中的位置。

对于一个像素点 $(u,v)$，程序先使用内参和畸变参数通过 \texttt{cv2.\allowbreak undistortPoints} 得到相机归一化射线，再将射线变换到世界坐标系，并与工作平面 $z=z_0$ 求交：
\[
P_W = O_W + \lambda d_W,\quad
\lambda = \frac{z_0-O_{W,z}}{d_{W,z}} .
\]
在视觉抓取实验中，还需要将棋盘格或相机工作平面坐标转换到机械臂基坐标。当前 \path{workspace_transform.json} 中采用的轴关系为：
\[
x_{\text{base}}=-y_{\text{cam}}+0.300,\quad
y_{\text{base}}=x_{\text{cam}}-0.028,\quad
z_{\text{base}}=z_{\text{cam}}.
\]
对应齐次矩阵为：
\[
T_{B}^{W}=
\begin{bmatrix}
0 & -1 & 0 & 0.300\\
1 & 0 & 0 & -0.028\\
0 & 0 & 1 & 0\\
0 & 0 & 0 & 1
\end{bmatrix}.
\]
这个矩阵表明，棋盘格平面坐标与机械臂基坐标之间存在一个平面内的 $90^\circ$ 轴向交换和一个平移偏置。后续视觉抓取中，目标像素中心先被反投影到棋盘格或工作平面坐标，再通过该矩阵转换到机械臂基坐标。

\begin{figure}[H]
  \centering
  \includegraphics[width=0.72\textwidth]{task4_camera_extrinsic_preview.jpg}
  \caption{棋盘格外参估计与坐标轴投影}
  \label{fig:task4-extrinsic}
\end{figure}

\subsection{实验步骤}

本实验代码使用已有棋盘格图像并剔除角点检测失败的样本，然后运行内参标定程序生成相机矩阵和畸变系数。获得内参后，程序利用 \texttt{solvePnP} 求解棋盘格坐标系相对于相机坐标系的位姿。最后，结合工作平面高度假设，将像素点反投影成空间射线，并求出其与桌面的交点。

实验过程中需要特别注意单位一致性。棋盘格方格边长以米为单位进入内参标定和外参估计，而机械臂控制程序中部分坐标采用毫米表示；如果单位混用，坐标转换结果会出现数量级错误。为避免这一问题，报告中统一在表格和公式中标注单位，并在代码分析中区分相机坐标、棋盘格坐标和机械臂基坐标。

为了检查外参是否合理，程序在图像上绘制棋盘格角点和三维坐标轴投影。若坐标轴方向与实际棋盘格摆放方向明显不一致，说明角点顺序、方格尺寸或外参求解存在问题；若坐标轴位置正确但重投影误差偏大，则可能是图像模糊、棋盘格没有铺平或内参不准确。图 \ref{fig:task4-extrinsic} 中坐标轴能够较好贴合棋盘格角点，因此可作为后续像素反投影的基础。

\subsection{标定结果与分析}

从内参标定结果看，平均重投影误差约为 0.20 px，说明棋盘格角点拟合精度较高。外参估计的重投影误差更低，约为 0.0215 px，表明在单幅棋盘格图像上估计得到的位姿与角点观测匹配良好。

但需要注意，图像重投影误差小并不等价于机械臂抓取误差一定小。抓取误差还会受到以下因素影响：

\begin{itemize}
  \item 相机安装位置和工作平面是否稳定；
  \item 棋盘格平面与真实物体上表面是否存在高度差；
  \item 像素中心提取是否受光照、遮挡和颜色反光影响；
  \item 机械臂自身运动学模型和舵机回差误差。
\end{itemize}

从机器人抓取角度看，标定误差会沿着“像素中心、相机射线、工作平面交点、机械臂基坐标”的链路逐级传递。即使像素级重投影误差很小，只要相机外参或桌面高度假设存在偏差，最终机械臂坐标也可能产生毫米级甚至厘米级误差。因此现有重投影误差只能说明图像角点拟合情况，不能证明真实机械臂抓取坐标准确。

本实验还说明了单目相机定位的一个基本限制：单个像素点只能确定一条空间射线，不能直接确定三维坐标。只有引入工作平面高度 $z=z_0$ 之后，射线与平面的交点才成为唯一的空间点。因此本实验适合桌面色块抓取这类目标高度近似固定的场景；若目标高度变化较大，或者物体堆叠，单一平面假设就会带来明显误差。此时需要引入深度相机、多视角重建或抓取前二次视觉确认。

\subsection{标定文件在系统中的作用}

相机标定实验最终生成多个文件，它们在后续视觉抓取中的作用不同。若只保存一张去畸变图像，无法支持实时抓取；若只保存内参而没有工作空间变换，也无法将像素点转换到机械臂坐标。因此标定结果需要按用途组织，如表 \ref{tab:calib-files} 所示。

\begin{table}[H]
  \centering
  \caption{相机标定输出文件及用途}
  \label{tab:calib-files}
  \small
  \begin{tabularx}{\textwidth}{L{4.0cm}Y Y}
    \toprule
    文件 & 主要内容 & 后续用途 \\
    \midrule
    \path{camera_calibration.json} & 图像尺寸、内参矩阵、畸变系数、重投影误差 & 去畸变、像素射线计算 \\
    \path{camera_calibration.npz} & NumPy 格式内参、畸变和新相机矩阵 & 供 Python 程序快速加载 \\
    \path{undistorted_example.jpg} & 去畸变效果图 & 检查标定结果是否合理 \\
    \path{camera_extrinsic.json} & 棋盘格到相机的外参、重投影误差 & 像素点投影和反投影 \\
    \path{camera_extrinsic_preview.jpg} & 坐标轴投影预览图 & 人工检查外参方向 \\
    \path{workspace_transform.json} & 工作平面到机械臂基坐标的轴关系和平移 & 将工作平面坐标转换为抓取坐标 \\
    \bottomrule
  \end{tabularx}
\end{table}

这些文件共同构成视觉抓取的坐标基础。实验五和实验六在识别算法上不同，但都必须依赖这里的标定结果。若相机位置移动、焦距变化、分辨率改变或棋盘格坐标系重新定义，相关标定文件都需要重新生成，否则后续抓取点会系统性偏移。由此可见，相机标定并不是一次性的“前处理”，而是视觉系统稳定运行的前置条件。

\subsection{小结}

本实验整理并运行了从像素坐标到工作平面坐标的标定计算程序。后续颜色识别和 YOLO 识别都需要利用相机内参、外参和工作空间到机械臂基坐标的转换关系生成候选抓取坐标，但该坐标尚未通过真实机械臂到点实验验证。
